\documentclass[12pt, titlepage]{article}

\usepackage{booktabs}
\usepackage{array}
\usepackage{longtable}
\usepackage{tabularx}
\usepackage{hyperref}
\hypersetup{
    colorlinks,
    citecolor=black,
    filecolor=black,
    linkcolor=red,
    urlcolor=blue
}
\usepackage[round]{natbib}

\input{../Comments}
\input{../Common}

\begin{document}

\title{Verification and Validation Report: \progname} 
\author{\authname}
\date{\today}
	
\maketitle

\pagenumbering{roman}

\section{Revision History}

\begin{tabularx}{\textwidth}{p{3cm}p{2cm}X}
\toprule {\bf Date} & {\bf Version} & {\bf Notes}\\
\midrule
Date 1 & 1.0 & Notes\\
Date 2 & 1.1 & Notes\\
\bottomrule
\end{tabularx}

~\newpage

\section{Symbols, Abbreviations and Acronyms}

\renewcommand{\arraystretch}{1.2}
\begin{tabular}{l l} 
  \toprule		
  \textbf{symbol} & \textbf{description}\\
  \midrule 
  T & Test\\
  \bottomrule
\end{tabular}\\

\wss{symbols, abbreviations or acronyms -- you can reference the SRS tables if needed}

\newpage

\tableofcontents

\listoftables %if appropriate

\listoffigures %if appropriate

\newpage

\pagenumbering{arabic}

This document ...

\section{System evaluation}
\subsection{Functional Requirements Evaluation}
\begin{longtable}{|>{\raggedright\arraybackslash}p{3cm}|>{\raggedright\arraybackslash}p{12cm}|}
\caption{Software Test Cases for User Account Creation} \\
\hline
\textbf{Test Case ID} & \textbf{user\_acc\_create} \\
\hline
\endfirsthead


\hline
Control & Automatic \\
\hline
Rationale & The system must be able to create and store a new account \\
\hline
State & User does not have an account on the application. \\
\hline
Input & Govt id (name, date of birth), license number (PT) OR invite code (patient), email, password (new) \\
\hline
Output & 'Account created successfully' displayed on page and user is successfully added to the user database \\
\hline
How the test will be performed & 
The test will be conducted with the use of a script which will input valid, invalid and edge case values \\
\hline

\textbf{Test Case ID} & \textbf{Invalid\_formatCreate} \\
\hline
Control & Automatic \\
\hline
Rationale & The data input by the user has to be in the correct format for our system to accept it as valid. \\
\hline
State & User does not have an account on the application. \\
\hline
Input & Incorrect format of license number/invite code OR incorrect format of date of birth, email, password(new) \\
\hline
Output & System displays 'incorrect format of one or more inputs' \\
\hline
How the test will be performed & 
The test will be conducted with the use of a script which will input valid, invalid and edge case values \\
\hline

\textbf{Test Case ID} & \textbf{Invalid\_sign\_up\_creds} \\
\hline
Control & Automatic \\
\hline
Rationale & Successful verification for both types of users is needed to use the application \\
\hline
State & User does not have an account on the application. \\
\hline
Input & Incorrect license number/mismatch of license number and govt id OR incorrect format of invite code, email, password(new) \\
\hline
Output & System displays 'License number cannot verified' OR 'invalid invite code' \\
\hline
How the test will be performed & 
The test will be conducted with the use of a script which will input valid, invalid and edge case values \\
\hline

\textbf{Test Case ID} & \textbf{Acc\_exists} \\
\hline
Control & Automatic \\
\hline
Rationale & New users should not already have an account on the system \\
\hline
State & Account already exists on system \\
\hline
Input & Govt id (name, date of birth), license number/invite code, email, password \\
\hline
Output & 'Account already exists. Wrong password. Try again.' \\
\hline
How the test will be performed & 
The test will be conducted with the use of a script which will input valid, invalid and edge case values \\
\hline

\end{longtable}



\section{Nonfunctional Requirements Evaluation}

\subsection{Usability}
		
\subsection{Performance}

\subsection{etc.}
	
\section{Comparison to Existing Implementation}	

This section will not be appropriate for every project.

\section{Unit Testing}

\section{Changes Due to Testing}

\wss{This section should highlight how feedback from the users and from 
the supervisor (when one exists) shaped the final product.  In particular 
the feedback from the Rev 0 demo to the supervisor (or to potential users) 
should be highlighted.}

\section{Automated Testing}
		
\section{Trace to Requirements}
		
\section{Trace to Modules}		

\section{Code Coverage Metrics}

\bibliographystyle{plainnat}
\bibliography{../../refs/References}

\newpage{}
\section*{Appendix --- Reflection}

The information in this section will be used to evaluate the team members on the
graduate attribute of Reflection.

\input{../Reflection.tex}

\begin{enumerate}
  \item What went well while writing this deliverable? 
  \item What pain points did you experience during this deliverable, and how
    did you resolve them?
  \item Which parts of this document stemmed from speaking to your client(s) or
  a proxy (e.g. your peers)? Which ones were not, and why?
  \item In what ways was the Verification and Validation (VnV) Plan different
  from the activities that were actually conducted for VnV?  If there were
  differences, what changes required the modification in the plan?  Why did
  these changes occur?  Would you be able to anticipate these changes in future
  projects?  If there weren't any differences, how was your team able to clearly
  predict a feasible amount of effort and the right tasks needed to build the
  evidence that demonstrates the required quality?  (It is expected that most
  teams will have had to deviate from their original VnV Plan.)
\end{enumerate}

\end{document}